\label{part:dynamic_process}

\textbf{(The Metatheory: Information Leadership \& Physical Support)}

\bluetext{双重立法的统一场}

\vspace{1em}长期以来，智能科学常被一条“笛卡尔裂痕”分成两类视角。第一类是\textbf{物理还原论 (Physical Reductionism)}，它把心智还原为神经电化学脉冲或逻辑门跳变。该视角具备物质实在性，但在解释“意义”与“目的”的高阶涌现时存在不足。第二类是\textbf{计算功能主义 (Computational Functionalism)}，它把智能抽象为脱离介质的算法与符号系统。该视角逻辑表达完备，但因弱化\textbf{热力学边界}而容易出现符号接地不足与系统“空心化”。

\vspace{1em}本部分作为本理论框架的公理化起点，目标是构建连接这两类视角的\textbf{拓扑桥梁}。为此提出\textbf{目的交互主义 (Purposeful Interactionism)}：智能过程具有不可分割的“波粒二象性”，既是物理世界中逆熵运行的\textbf{耗散结构 (Dissipative Structure)}，也是信息生态中定义价值的\textbf{语义主体 (Semantic Subject)}。

\vspace{1em}在此，我们将确立智能演化的\textbf{“双层立法” (Dual Legislation)} 机制，这构成了全书核心拉格朗日量 $\mathcal{L}_{total}$ 的两大支柱：
\begin{enumerate}
        \item \textbf{信息主导 (Information Leadership)}：作为\textbf{立法者 (The Legislator)}。
        \begin{itemize}
                \item 它源于主客体的交互，定义了系统的\textbf{“几何扩张倾向”}。
                \item 它通过\textbf{体验图 ($G_E$)} 定义价值势能，告诉几何流形“应当如何弯曲”以最大化语义收益。它是\textbf{目的论的拉力}。
        \end{itemize}
        \item \textbf{物理支撑 (Physical Support)}：作为\textbf{执法者 (The Enforcer)}。
        \begin{itemize}
                \item 它源于热力学与动力学定律，定义了系统的\textbf{“能量收缩约束”}。
                \item 它通过\textbf{兰道尔极限}与\textbf{介质粘滞}定义能耗边界，告诉几何流形“只能如何演化”以最小化物理代价。它是\textbf{因果论的阻力}。
        \end{itemize}
\end{enumerate}
智能的动力学特征可视为这两股力量在\textbf{变分博弈}中的解：系统在物理法则给定的\textbf{边界条件}下，为实现信息价值最大化而进行\textbf{时空几何重构 (Spacetime Geometric Reconstruction)}。

\vspace{1em}\redtext{前文已经讨论了智能体在环境信息表示上的静态机制（MSC）。从本部分开始，我们在 MSC 基础上讨论智能的动态生成过程（HSF-HD），并据此推导相关第一性原理。}



%---chp----
